%  LaTeX support: latex@mdpi.com
%=================================================================
%  Built against the official MDPI template.
%
%  REQUIRED FILES: Definitions/ from the MDPI template package, plus
%  fig_validation.pdf, fig1_ceiling_definition.pdf and
%  fig_inversion_wide.pdf in the same directory.
%
%  SUGGESTED SPECIAL ISSUE (name it on the submission form):
%      "Thermal and Optical Remote Sensing: Evaluating Urban Green
%       Spaces and Urban Heat Islands in a Changing Climate"
%       https://www.mdpi.com/journal/remotesensing/special_issues/RS_Urban_Green_UHI
%  DEADLINE NOT VERIFIED -- mdpi.com blocks automated fetching.
%
%  BEFORE SUBMITTING:
%    * Nine reference entries still carry "et al."; MDPI wants full
%      author lists.
%    * Verify [friedman2001greedy] volume and page range.
%    * Make the GenAI disclosure in cknowledgments match what you
%      actually did. MDPI requires it.
%=================================================================

\documentclass[remotesensing,article,submit,oneauthor]{Definitions/mdpi}

\firstpage{1}
\makeatletter
\setcounter{page}{\@firstpage}
\makeatother
\pubvolume{1}
\issuenum{1}
\articlenumber{0}
\pubyear{2026}
\copyrightyear{2026}
\datereceived{ }
\daterevised{ }
\dateaccepted{ }
\datepublished{ }

\Title{Data Ceilings in Urban Heat Island Modeling Are Conditional on
Spatial Support: A Land-Cover Diagnostic Applied Across Five Continents}

\Author{Abhiram Donepudi}

\AuthorNames{Abhiram Donepudi}

\address{%
$^{1}$ \quad Westwood High School, Round Rock, TX 78681, USA;
abhiram.donepudi@gmail.com}

\corres{Correspondence: abhiram.donepudi@gmail.com}

\abstract{When an urban heat island (UHI) model transfers poorly to a new
city, there are two explanations. The model may have failed, or the land
cover may not have separated thermal regimes well enough for a
land-cover model to succeed. We introduce a land-cover data ceiling that
separates these issues. It is independent of the simulator, requiring a
city boundary, a land cover raster, and one satellite scene. We define the
ceiling as the multiple correlation between MODIS land surface temperature
and the land-cover channels a model uses. Impervious fraction alone is
invalid, as it goes negative in arid cities even where the ceiling exceeds
$R = +0.45$. A ceiling constrains a model only when their receptive fields
match. In our Austin test, the simulator exceeds the per-cell ceiling
on held-out days, while calculating the ceiling over the model's
extent restores the bound. Built--vegetated land-cover difference predicts
the ceiling across 49 US cities (Spearman $\rho = +0.69$, $p = 3 \times
10^{-8}$). Across 28 cities on five continents, all eleven inverting cities
recover to ceilings between $+0.47$ and $+0.91$ once neighborhoods are
discernible. Among cities we could observe, none is separability-limited,
and the arid/non-arid difference is no longer detectable at matched
receptive field ($p = 0.115$), versus $p = 0.018$ at pixel scale.}

\keyword{urban heat island; land surface temperature; MODIS; class
separability; model transferability; surface energy balance; open data;
thermal remote sensing}

\begin{document}

\section{Introduction}

A common phenomenon in UHI model validation is conflating two failure
modes when translating a model to another city. The model may have
actually failed to capture that city's thermal structure, or the input
land-cover data may not separate thermal regimes well enough for
\emph{any} land-cover-driven model to succeed. Those demand opposite
solutions: the first needs a more capable model, and the second needs
better inputs, or an admission that the city can't be represented with
these channels. But both result in the same outcome, poor correlation
against observed LST, and the literature offers no standard way to tell
them apart.

This paper introduces a \emph{land-cover data-ceiling control} that makes
the distinction explicit. The ceiling is a per-city bound on the skill
recoverable from a given set of land-cover inputs, computed from
observations alone. Because it involves no simulator, calibration or
meteorological forcing, it is not a property of the model under test.

A ceiling is only a bound relative to a stated model class, and we show
that the class must be matched in two respects rather than one. The
channel count is the obvious requirement. The less obvious and more
consequential one is the \emph{receptive field}, as a ceiling regressed
on each cell's own land cover does not bound a model that reads across
space, and our own simulator exceeds such a ceiling in its reference
city. Matching the ceiling's spatial reach to the model's reach restores
the bound, and changes which cities are judged
separability-limited.

We exercise the diagnostic on \textsc{austin-twin}, a physics-based slab
energy-balance simulator driven entirely by zero-authentication open data
\cite{oke1982energetic,voogt2003thermal,masson2000physically}. The
simulator is an experimental instrument here, not the scientific
endpoint. We are not reporting the performance of a heat model; we are
using a model whose behaviour we control to show what a diagnostic does
and does not bound. Its simplifications are therefore a feature of the
design---a model with fully realistic physics would make it harder, not
easier, to attribute a ceiling violation to receptive field rather than to
something else. Existing UHI tools span two
extremes. High-fidelity mesoscale models, such as WRF-SLUCM, TEB, and
the Urban Weather Generator \cite{bueno2013urban,kusaka2001simple},
resolve canopy processes but require proprietary or authenticated
inputs and hours of compute per scenario. Machine-learning approaches
\cite{lyu2022integrated,hoang2025geospatial,bhamjee2024detection}
validate well against historical observations but cannot predict the
effect of an intervention outside their training distribution, which
is what urban planning asks of them. \textsc{austin-twin} occupies the
middle, as it is a 2D grid model running in 10.8~ms per 48-hour
simulation on a desktop CPU, with an explicit counterfactual engine and
fully open inputs. Table~\ref{tab:novelty} positions it. The data-ceiling
column has no prior entry.

Another methodological motivation concerns reproducibility. Most urban
climate pipelines depend on authenticated inputs, including reanalysis
behind a registration wall, licensed imagery, and institutional data
agreements. However, that dependency is invisible in a published result
and harmful to independent verification, because a reader who cannot
obtain the inputs cannot check the outputs. Here, every input used is
retrievable by anonymous HTTP request.

The central claim is not that we have measured a data ceiling. It is
that \emph{any} empirical data-ceiling diagnostic is conditional on
spatial support, and that using a mismatched support can systematically
misdiagnose model failure as data limitation. The object we estimate is
accordingly not the maximum correlation the data permit; it is an
empirically estimated benchmark conditioned on an explicitly specified
information set and spatial support, and we report it as such throughout.

Around that claim sit five reusable results, each independent of this
simulator:

\begin{enumerate}
\item A \emph{land-cover data-ceiling control}, together with the two
specification requirements that make it a bound rather than a number.
It must be defined across all channels a model consumes---the univariate
alternative goes negative in arid cities---and its receptive field must
match the model's, since a per-cell ceiling does not bound a model that
reads across space. We pair it with a built--vegetated thermal-gap
heuristic validated across 49 cities, and an NDVI channel screen shown
to be a predictable, general rule rather than a single-city result.

\item A demonstration that receptive field is not a technicality: our
own simulator exceeds Austin's per-cell ceiling on its held-out days,
and matching the ceiling's spatial reach both restores the bound and
reverses which of two transfer-failure cities is judged
separability-limited.

\item Evidence that the arid inversion of the land-cover--temperature
relationship is \emph{climatic} rather than continental, established on
28 cities across five continents under a city definition uniform enough
to support the comparison---together with the finding that the effect is
definition-sensitive in the American Southwest, which we report rather
than omit.

\item A \emph{decay-constrained two-stage calibration} quantifying the
trade-off between MODIS pattern fidelity and intervention spatial
realism.

\item An \emph{evapotranspiration ablation} showing that a standard
linear evapotranspiration (ET) term contributes nothing to
spatial-pattern skill at 500~m, regardless of formulation.
\end{enumerate}

The last two are reported because a ceiling diagnostic is only as
trustworthy as the model it is demonstrated on. A reader is entitled to
ask whether the simulator's transfer failures reflect the diagnostic or
merely a weak simulator, and the diffusion and ET results establish what
the vehicle can and cannot do. They are supporting evidence for the first
three rather than independent contributions.


\begin{table}[H]
\caption{Positioning relative to existing UHI simulation and prediction
systems. ``Zero-auth'' means all inputs are retrievable by anonymous HTTP
request; ``counterfactual'' means the system evaluates an intervention
outside the training distribution via an explicit causal mechanism;
``ceiling diagnostic'' means a data-ceiling control separating model
failure from input class-separability limits.\label{tab:novelty}}
\begin{tabularx}{\textwidth}{lCCCCC}
\toprule
\textbf{System} & \textbf{Scale} & \textbf{Zero-Auth} &
\textbf{Counterfactual} & \textbf{Sat.\ Calib.} & \textbf{Ceiling Diag.} \\
\midrule
WRF-SLUCM/TEB \cite{kusaka2001simple,masson2000physically} & Meso/canyon & No & Yes & Rare & No \\
Urban Weather Generator \cite{bueno2013urban} & Neighborhood & No & Yes & No & No \\
ML (NGBoost/CNN/GeoFM) \cite{lyu2022integrated,hoang2025geospatial} & City--regional & Partial & No & Implicit & No \\
HeatMat \cite{heatmat2026} & District & Partial & No & Yes & No \\
\textbf{This work} & City & Yes & Yes & Yes & Yes \\
\bottomrule
\end{tabularx}
\end{table}

\section{Materials and Methods}

\subsection{Study Areas and City Definitions}

Austin, Texas is the reference city, with Phoenix (AZ), Miami (FL) and
Denver (CO) as transfer targets chosen to span the humid--arid range. The
wide sweep covers 49 US cities (55 attempted, 6 dropped for insufficient
cloud-free coverage). The global extension covers 28 cities on five
continents.

City delimitation requires care once the sample leaves the United States.
OpenStreetMap administrative units are not internationally comparable:
querying ``Riyadh, Saudi Arabia'' returns a 375,478~km$^2$ province,
``Perth'' a 6~km$^2$ council area, and ``Alice Springs'' 1~km$^2$. All
three return a valid geometry and would produce a number that is not reflective of the true value. For the global extension, therefore, every city (US and non-US alike) is defined as a 25~km-radius disc about its center (1960~km$^2$, within the
US sample's own range of administrative areas). Including US cities under
this identical definition provides a comparability bridge, so that any
US/non-US difference cannot be an artifact of how a city was delimited.
Each city is sampled in its own warm season, with southern-hemisphere
cities drawing on an austral-summer date pool.

\subsection{Data}

All inputs are retrieved anonymously, without API keys, accounts or
licenses. Austin's administrative boundary comes from OpenStreetMap.
Land cover uses ESA WorldCover v200 at 10~m \cite{zanaga2022esa},
aggregated to a 500~m UTM grid ($\sim$2896 cells for Austin) as class
fractions rather than a majority class, giving three channels: impervious,
vegetation and water fraction. MODIS Aqua MYD11A1 v6.1 daily LST is
retrieved through the Microsoft Planetary Computer STAC API
\cite{microsoft_pc}. The Aqua afternoon overpass ($\sim$13:30 local) was
deliberately chosen to be used, as it falls near peak surface temperature where the UHI signal is strongest. LST is read in its native scaled-integer form with
the fill value masked before conversion, and reprojected by
nearest-neighbor resampling so that fill values cannot bleed into valid
neighbors.

MYD11A1 is natively 1~km, so the 500~m analysis grid oversamples it and
each thermal pixel spans roughly four cells. Nearest-neighbor resampling
is deliberate---it preserves the fill value, which averaging would
smear---but it means the effective number of independent thermal
observations is well below the cell count $n$ we report, by a factor of
about four from resampling and further where the thermal field is
genuinely homogeneous. Two consequences are handled explicitly: all
uncertainty estimates use spatial blocks wider than the thermal pixel
(Section~\ref{sec:stats}), and the receptive-field result is re-tested on
coarser supports where predictor and target share a pixel
(Section~\ref{sec:support}). ERA5 hourly reanalysis \cite{hersbach2020era5} is read from
Google's analysis-ready cloud-optimized (ARCO) Zarr mirror, which is
anonymously readable. This replaces the registration-gated Climate Data
Store API. MODIS NDVI (MOD13A1, 500~m) supports the channel screen.

\subsection{Surface Energy Balance Simulator}

Each grid cell evolves under a slab energy balance,
\begin{equation}
\frac{dT_i}{dt} = Q_{\mathrm{solar}} - Q_{\mathrm{rad}}
                - Q_{\mathrm{ET}} + Q_{\mathrm{diff}},
\label{eq:slab}
\end{equation}
with per-class solar absorption, Newtonian radiative cooling, linear
evapotranspiration $Q_{\mathrm{ET}} = k_{\mathrm{ET}} f_{\mathrm{veg}}
\max(T_i - T_{\mathrm{air}}, 0)$, and a Fickian diffusion term. Water
cells use an explicit mixed-layer heat capacity and a fixed latent-heat
flux rather than a single relaxation parameter. The timestep satisfies
the Courant--Friedrichs--Lewy condition for every diffusivity in the
search range.

\subsection{Calibration}

Five coefficients are tuned by differential evolution
\cite{storn1997differential} against three training MODIS days, with two
held out for testing. The \emph{ship configuration} uses a two-stage
procedure: fix $D = 1.0$~m$^2$~s$^{-1}$, matching the literature
tree-cooling half-distance \cite{ziter2019scale}, then refit the
remaining four coefficients under ARCO-ERA5 forcing. Ship coefficients are
transferred to Phoenix, Miami and Denver under three schemes: zero-shot;
Scheme~B, re-optimizing $k_{\mathrm{ET}}$ only; and Scheme~A,
re-optimizing all four coefficients. All configurations being compared use
identical optimizer settings and seed, so that differences between them
are not confounded with differences in search effort.

\subsection{The Land-Cover Data Ceiling}

A \emph{data-ceiling control} estimates the correlation achievable in a
given city from the information the inputs carry. It is not an
unconditional upper bound on prediction, and we write it so that its
conditioning is explicit:
\begin{equation}
C\!\left(X, Y; R\right),
\label{eq:ceiling}
\end{equation}
the achievable correlation between target $Y$ and channel set $X$ for a
model whose receptive field is $R$. All three arguments matter. Changing
the channels changes it, changing the target changes it, and---the result
this paper is mainly about---changing $R$ changes it, by more than either
of the others in the cities where it matters most.

Quoting a percent-of-ceiling figure therefore requires naming the model
class it is conditioned on. A model whose receptive field exceeds $R$ can
exceed $C(X,Y;R)$ without anything being wrong, and we demonstrate exactly
that in Section~\ref{sec:cnn}. We retain the word ceiling because the
quantity does what a ceiling is wanted for---it separates information
limits from model limits---but it ceilings a \emph{stated} model class,
not prediction in general.

We also distinguish two limitations that a single phrase would conflate. A
city is \emph{separability-limited} when its land cover cannot distinguish
thermal regimes at the scale being read, which is what $C(X,Y;R)$
measures. A city is \emph{observability-limited} when the thermal
observations needed to compute anything at all do not exist, which
$C(X,Y;R)$ is silent about by construction. Both are data limits; only the
first is a statement about land cover. We use the specific terms
throughout and reserve ``data-limited'' for cases where we mean both.

An achievable-performance bound conditioned on the available inputs is not
itself a new idea. In classification it is the Bayes error, the
irreducible error of the optimal decision rule, and Fano's inequality
relates it to the mutual information between features and target
\cite{cover2006elements}. Estimating it from finite data without first fitting an
optimal model is an active problem in its own right \cite{noshad2019learning}. Our
$C(X,Y;R)$ is a regression-side analogue of the same quantity, and we make
no claim to have invented the notion of a data-imposed limit.

Two things separate what we report from that literature. The first is the
conditioning itself: Bayes error is defined relative to a feature set,
whereas $C(X,Y;R)$ is defined relative to a feature set \emph{and} a
spatial support, and our results show the second argument moves the
quantity more than the first in the cities where it matters. Spatial
support has no analogue in the classical setting, where observations are
exchangeable. The second is practical. We attempted a distribution-free
estimate along the classical lines, using a $k$-nearest-neighbour mutual
information estimator converted to a correlation equivalent, and abandoned
it: on Miami it returned $+0.58$ against an achievable $+0.89$, violating
the bound it was supposed to provide, because the Gaussian inversion fails
on a bimodal land-water thermal distribution. At the sample sizes a single
city affords, an empirical achievable-skill estimate proved more
trustworthy than an information-theoretic bound, which is why we report
one and say plainly that it is a lower bound.


\subsubsection{Channel Count}

The first is straightforward. A model consuming three channels can exceed
a one-channel correlation by construction, so a ``percent of ceiling
achieved'' computed against the univariate $r(\mathrm{LST},
\text{impervious})$ is not interpretable. Worse, in arid cities the
univariate quantity goes negative even where the three-channel fit exceeds
$R = +0.45$, because desert vegetation runs hotter than built surfaces
there, leading to a negative denominator, making the ratio meaningless. We therefore define the ceiling over all channels the model consumes, as a multiple
correlation $R$.

\subsubsection{Receptive Field}
\label{sec:receptive}

The second is less obvious and, we find, more consequential. A multiple
correlation regresses LST on each cell's \emph{own} land cover. A model
with any spatial reach does not work that way: our simulator diffuses
heat, so a modeled cell depends on its neighbors' land cover as well.
Convolutional and foundation-model approaches
\cite{lyu2022integrated,hoang2025geospatial} have receptive fields wider
still. Such a model has access to information a per-cell ceiling never
offered its estimator, and can exceed that ceiling with nothing wrong.

This is not hypothetical, as our own shipped configuration exceeds Austin's
per-cell ceiling on the very days it is scored on
(Section~\ref{sec:rfresult}), which is what led us to the distinction.

We therefore also report a \emph{receptive-field-matched} ceiling,
computed by giving the estimator mask-aware neighborhood averages of the
same three channels at $\sigma = 1, 2, 4$ cells (roughly 0.5--2~km),
bracketing the shipped diffusion length of $\approx$590~m. Smoothing is
applied to land cover, which is a known input everywhere in the city, and
not to LST, so no label information is introduced.

\subsubsection{Estimation}

Both ceilings are estimated two ways: as a multiple correlation, and as
the out-of-fold correlation of a gradient-boosted regressor
\cite{friedman2001greedy}, which removes the linearity assumption. The
latter uses \emph{spatial} block cross-validation over contiguous
$6 \times 6$-cell tiles.

We call this quantity a ceiling for continuity with the linear case, but
it should be read precisely: it is the out-of-fold skill of one benchmark
regressor, so it is a \emph{lower bound} on what the inputs can support,
not a proven upper bound on all models. A better estimator, or a wider
receptive field, could raise it. That asymmetry is why we treat a model
exceeding a ceiling as evidence the ceiling is mis-specified rather than
as evidence about the model.

The blocking is not optional. Adjacent 500~m cells are
strongly autocorrelated, so random folds train on a test cell's own
neighbors and the score inflates toward unity, fabricating the
conclusion that land cover explains everything.

Because single-day ceilings are noisy, we report multi-day means
throughout (5--12 cloud-free dates per focus city; median 6 across the
49-city sample).

\subsection{Statistical Methods}
\label{sec:stats}

Model skill is reported as the spatial correlation $r$ and RMSE of
\emph{mean-subtracted} fields. The slab formulation carries no
absolute-temperature anchor, so only the spatial pattern is meaningfully
scored; an RMSE computed on raw fields additionally carries a constant
offset the model does not attempt to predict, and would understate
agreement for reasons unrelated to pattern skill. We state this
explicitly because recomputing RMSE the obvious way yields a substantially
different number.

Because neighboring grid cells are strongly autocorrelated, per-city
confidence intervals come from a spatial block bootstrap resampling
contiguous $6 \times 6$-cell tiles ($B = 2000$); a cell-wise bootstrap
would understate intervals by treating each cell as independent.
Comparisons between two configurations evaluated on the same cities use a
\emph{paired} bootstrap, so the interval describes the difference rather
than being inferred from two overlapping marginal intervals.

City-level associations, such as gap against ceiling, and aridity against
ceiling, use Spearman rank correlation, since neither relationship is
assumed linear and single extreme cities should not dominate. Group
comparisons (arid versus humid; arid US versus arid non-US) use the
Mann--Whitney $U$ test, reported with group sizes, because these samples
are small enough that a non-significant result may reflect limited power
rather than the absence of an effect.

We report all tests uncorrected for multiple comparisons. The study is
structured around a small number of pre-specified comparisons---gap
against ceiling, arid against non-arid, per-cell against
receptive-field-matched---and the remaining tests are exploratory. Readers
should treat individual $p$-values between roughly $0.01$ and $0.1$
accordingly; we flag the two that fall in that range, the humid subset and
the matched arid/non-arid contrast, wherever they appear.

\section{Results}
\label{sec:results}

\subsection{Calibration and Uncertainty}

Table~\ref{tab:bootstrap_ci} reports spatial block-bootstrap 95\%
confidence intervals. The ship configuration achieves $r = +0.754$
[$+0.718$, $+0.785$], trading 0.10 in $r$ against an unconstrained
single-stage fit ($r = +0.850$) in exchange for physically realistic
diffusion and intervention decay lengths (Section~\ref{sec:disc}). ERA5
forcing contributes $+0.007$ [$+0.004$, $+0.009$] over synthetic
forcing, which is small yet distinguishable from zero, with the benefit being
mainly methodological, in that it permits absolute-temperature
comparability without anomaly normalization.

\begin{table}[H]
\caption{Block-bootstrap 95\% confidence intervals, held-out test $r$ and
RMSE. Both statistics are computed on mean-subtracted
fields.\label{tab:bootstrap_ci}}
\begin{tabularx}{\textwidth}{lCC}
\toprule
\textbf{Configuration} & \textbf{Test} {\boldmath$r$} \textbf{[95\% CI]} & \textbf{RMSE (}{\boldmath$^{\circ}$}\textbf{C)} \\
\midrule
Single-stage, synthetic forcing & $+0.850$ [$+0.823$, $+0.872$] & 1.09 \\
Two-stage, synthetic forcing & $+0.747$ [$+0.712$, $+0.778$] & 1.48 \\
\textbf{Ship (ARCO-ERA5)} & {\boldmath$+0.754$} [$+0.718$, $+0.785$] & \textbf{1.40} \\
\bottomrule
\end{tabularx}
\end{table}

Figure~\ref{fig:validation} shows the simulated field against the MODIS
observation for a held-out day, in both absolute and anomaly form. The
model reproduces the hot urban core and the cool riparian corridor, but
runs a mean 3.66~$^{\circ}$C cold in absolute terms (panel~c). The
anomaly panels (d--f) show what the reported statistics actually score.
The residual is not spatially uniform, a point we return to in
Section~\ref{sec:disc}.

\begin{figure}[H]
\includegraphics[width=\textwidth]{fig_validation.pdf}
\caption{Simulated versus observed LST for Austin, 21 August 2024, a
held-out day. (\textbf{a}) Simulated peak-heat field under the ship
configuration; (\textbf{b}) MODIS Aqua observation; (\textbf{c}) raw
difference, which carries the $-3.66$~$^{\circ}$C mean offset;
(\textbf{d},\textbf{e}) the same two fields after mean subtraction;
(\textbf{f}) residual of the anomalies, which is what the reported
$r = +0.739$ and anomaly RMSE of 1.49~$^{\circ}$C are computed from
($n = 2868$ cells). Panels (a,b) share one color scale, as do
(d,e).\label{fig:validation}}
\end{figure}
\unskip

\subsection{The Ceiling Must Match the Model's Receptive Field}
\label{sec:rfresult}

A per-cell ceiling does not bound a spatial model. The clearest
demonstration of this is our own. On the two held-out days of
Table~\ref{tab:bootstrap_ci}, the shipped configuration scores
$r = +0.754$ while Austin's per-cell multivariate ceiling on those same
days is $+0.693$. The model stands at 109\% of its own ceiling.

Relaxing linearity does not account for this. A gradient-boosted
regressor on the same three per-cell channels, under spatial block
cross-validation, reaches $+0.681$ on those days, which is marginally
\emph{worse} than the linear fit, and still exceeded by the simulator.
Across the four focus cities the per-cell nonlinear gain is small and
inconsistently signed (Austin $-0.024$, Denver $-0.014$, Miami $+0.047$,
Phoenix $+0.110$), so whatever the simulator is exploiting, it is not
nonlinearity.

Instead, it is spatial reach. Giving the estimator the neighborhood-averaged
channels of Section~\ref{sec:receptive} raises Austin's ceiling on those
days to $+0.883$, placing the model at 85\% of it, so the violation
disappears. Table~\ref{tab:receptive} gives the multi-day picture.

\begin{table}[H]
\caption{Data ceilings under three specifications, multi-day means. The
context column adds mask-aware neighborhood averages of the same three
channels at $\sigma = 1, 2, 4$ cells. Gain is the context ceiling minus
the per-cell nonlinear one.\label{tab:receptive}}
\begin{tabularx}{\textwidth}{lCCCC}
\toprule
\textbf{City} & \textbf{Linear, per-Cell} & \textbf{Nonlinear, per-Cell} &
\textbf{Nonlinear, Context} & \textbf{Gain} \\
\midrule
Austin & $+0.613$ & $+0.589$ & $+0.789$ & $+0.200$ \\
Phoenix & $+0.271$ & $+0.381$ & $+0.689$ & $+0.307$ \\
Denver & $+0.437$ & $+0.423$ & $+0.756$ & $+0.333$ \\
Miami & $+0.841$ & $+0.887$ & $+0.952$ & $+0.065$ \\
\bottomrule
\end{tabularx}
\end{table}

The gain is not an artifact of smoothing crossing fold boundaries. At
block sides of 6, 12 and 24 cells, the last four times the smoothing
reach, Austin gains $+0.202$, $+0.208$ and $+0.170$, and Phoenix
$+0.310$, $+0.252$ and $+0.290$.

Two consequences follow. Firstly, a percent-of-ceiling figure is only
interpretable once the ceiling's receptive field matches the model's;
quoted against a per-cell ceiling, it can exceed 100\% for any model with
spatial reach. Second, and more practically, the gain is largest in the
arid cities and smallest in Miami, which we return to in
Section~\ref{sec:disc}.

\subsubsection{The Gain Is Not an Artifact of Thermal Resolution}
\label{sec:support}

Because MYD11A1 is natively 1~km while land cover varies at 500~m, the
predictors sit at a finer support than the target. Smoothing land cover at
$\sigma = 1$--$2$ cells is close to what is needed to bring the two into
alignment, so the gain reported above admits a second explanation:
\emph{support matching}, an ordinary geostatistical effect, rather than
receptive field. At 500~m both explanations predict the same observation.

Coarser supports separate them. We block-aggregated land cover
\emph{and} LST onto a common support at 1000~m and 1500~m, keeping the
smoothing scales strictly wider than the new pixel, so that a surviving
gain cannot be produced by support matching (Table~\ref{tab:support}).

\begin{table}[H]
\caption{Context gain under progressively coarser common support. Land
cover and LST are block-aggregated together, and smoothing operates only
beyond the resulting pixel. Denver has too few cells to score at
1500~m.\label{tab:support}}
\begin{tabularx}{\textwidth}{lCCC}
\toprule
\textbf{City} & \textbf{500 m} & \textbf{1000 m} & \textbf{1500 m} \\
\midrule
Austin & $+0.200$ & $+0.151$ & $+0.094$ \\
Phoenix & $+0.307$ & $+0.310$ & $+0.294$ \\
Denver & $+0.333$ & $+0.383$ & --- \\
Miami & $+0.065$ & $+0.048$ & $+0.037$ \\
\bottomrule
\end{tabularx}
\end{table}

The result splits by city, and informatively. In the two low-per-cell
ceiling cities the gain is stable or grows: Phoenix holds at $+0.307$,
$+0.310$, $+0.294$ and Denver rises from $+0.333$ to $+0.383$. These are
the cities the finding is about, and their gain persists where smoothing
acts beyond the thermal pixel, which support matching cannot produce.

In the two high-ceiling cities it decays: Austin more than halves from
$+0.200$ to $+0.094$, Miami from $+0.065$ to $+0.037$. There, a
substantial share of the 500~m gain \emph{was} support matching, and the
500~m figures overstate the effect. We report the decay rather than only
the survival because it bounds the claim: receptive field is what recovers
signal in cities whose per-cell ceilings are low, and the correction is
smaller than the 500~m numbers suggest wherever the per-cell ceiling was
already high.

\subsubsection{A Convolutional Model Behaves as the Argument Predicts}
\label{sec:cnn}

The argument so far rests on a diffusion simulator, yet its stated
relevance is to convolutional and foundation-model approaches, whose
receptive fields are wider still. We therefore tested it directly. A
fully-convolutional network was trained to predict LST from the same
three land-cover channels, under the same spatial block folds, at three
receptive fields: 1 cell (all $1{\times}1$ convolutions, a per-cell model
that is still nonlinear), 7 cells (three stacked $3{\times}3$ layers,
$\approx$3.5~km) and 13 cells (six layers, $\approx$6.5~km). Only the LST
target is masked out of the loss; land cover is an input available
everywhere in the city, so no label information crosses a fold boundary.

\begin{table}[H]
\caption{Convolutional models against the two ceilings, multi-day means.
RF~=~1 is the control: a neural network with no spatial reach. Bold marks
figures that exceed the context ceiling.\label{tab:cnn}}
\begin{tabularx}{\textwidth}{lCCCCC}
\toprule
\textbf{City} & \textbf{Per-Cell Ceiling} & \textbf{CNN, RF = 1} &
\textbf{Context Ceiling} & \textbf{CNN, RF = 7} & \textbf{CNN, RF = 13} \\
\midrule
Austin & $+0.589$ & $+0.605$ & $+0.789$ & $+0.785$ & {\boldmath$+0.862$} \\
Phoenix & $+0.381$ & $+0.407$ & $+0.689$ & $+0.608$ & {\boldmath$+0.767$} \\
Denver & $+0.423$ & $+0.492$ & $+0.756$ & $+0.693$ & {\boldmath$+0.842$} \\
Miami & $+0.887$ & $+0.876$ & $+0.952$ & $+0.949$ & {\boldmath$+0.956$} \\
\bottomrule
\end{tabularx}
\end{table}

Three things follow, and the first is the control.

\textbf{Architecture alone buys nothing.} The RF~=~1 network lands within
$0.07$ of the per-cell ceiling in every city, and \emph{below} it in
Miami. A neural network denied spatial reach does not beat gradient
boosting on the same per-cell inputs. Whatever produces the gain, it is
not the model family.

\textbf{Spatial reach buys a great deal.} RF~=~7 exceeds the per-cell
ceiling by $+0.196$, $+0.227$, $+0.270$ and $+0.062$ respectively---the
same ordering across cities as the Gaussian-context gain of
Table~\ref{tab:receptive}, which is what one would expect if the two are
measuring the same thing.

\textbf{The matched ceiling bounds a matched model.} RF~=~7 sits at or
below the context ceiling in all four cities, and in two of them almost
exactly at it: $+0.785$ against $+0.789$ in Austin, $+0.949$ against
$+0.952$ in Miami. This is direct evidence that the receptive-field-matched
ceiling functions as a bound, rather than merely being a larger number
than the per-cell one.

\textbf{And a wider model breaks it.} RF~=~13 exceeds the context ceiling
in all four cities, by $+0.073$, $+0.078$, $+0.086$ and $+0.004$.

We read this as the experiment's most informative outcome rather than as a
discrepancy to be explained away. What these results establish is that an
achievable-correlation estimate is \emph{necessarily} conditional on
spatial support: a model reading $\approx$6.5~km has access to structure
that a three-Gaussian specification bracketing a 590~m diffusion length
never offered its estimator, and it uses it. What they do not establish,
and what we do not claim, is that our particular specification is a
universal upper bound. The general statement is the defensible one and the
stronger one. A ceiling is not a property of a city; it is
$C(X,Y;R)$, and $R$ must be named before a percent-of-ceiling figure means
anything.

\subsection{Multi-City Transfer and the Ceiling Diagnostic}

Table~\ref{tab:transfer} presents the transfer matrix. Against per-cell
ceilings, Miami reaches 99\% under Scheme~A ($r = +0.832$ against
$+0.841 \pm 0.025$), while Phoenix and Denver reach 28\% and 29\%
($+0.271 \pm 0.085$ and $+0.437 \pm 0.085$).

However, read against receptive-field-matched ceilings, the picture
changes materially. Phoenix and Denver have ceilings of $+0.689$ and
$+0.756$ (Table~\ref{tab:receptive}), so the achieved fractions become
11\% and 17\%. Both cities retain substantial unrecovered headroom, and
the earlier reading (that Denver's higher ceiling leaves room for
calibration to exploit while Phoenix's does not) is not supported.
Re-optimization still degrades Phoenix ($+0.077 \rightarrow +0.033$) and
roughly doubles Denver ($+0.056 \rightarrow +0.126$), but the absence of
recoverable signal in Phoenix is no longer the explanation: the signal is
present, and this simulator is not capturing it.

We report this interpretation rather than the more convenient earlier framing because it is what the matched ceilings demonstrate and clarifies the diagnostic's purpose: the ceiling bounds what is recoverable, but does not diagnose \emph{why} a particular model fails to recover it.

\begin{table}[H]
\caption{Multi-city transfer, held-out test $r$ over 2 days. The ceiling
is the multi-day mean $R$(LST; impervious, vegetation,
water).\label{tab:transfer}}
\begin{tabularx}{\textwidth}{lCCC}
\toprule
\textbf{Scheme} & \textbf{Denver} & \textbf{Phoenix} & \textbf{Miami} \\
\midrule
Austin ship (reference) & $+0.754$ & $+0.754$ & $+0.754$ \\
Zero-shot & $+0.056$ & $+0.077$ & $+0.724$ \\
Scheme B, 3 days & $+0.063$ & $+0.071$ & $+0.776$ \\
Scheme A, 3 days & $+0.126$ & $+0.033$ & $+0.832$ \\
\midrule
\textbf{Ceiling} & {\boldmath$+0.437 \pm 0.085$} & {\boldmath$+0.271 \pm 0.085$} & {\boldmath$+0.841 \pm 0.025$} \\
\bottomrule
\end{tabularx}
\end{table}

Figure~\ref{fig:ceiling} shows why the univariate definition fails. Every
city lies on or above the $y = x$ line, because the multiple correlation
is greater than or equal to the univariate one by construction. A
three-channel model can therefore exceed a one-channel ``ceiling'' with
nothing wrong. Five cities fall left of $x = 0$, where a percentage of
ceiling is undefined.

\begin{figure}[H]
\includegraphics[width=\textwidth]{fig1_ceiling_definition.pdf}
\caption{Univariate against multivariate data ceiling for the 49-city US
sample. Points on or above the diagonal show that the multiple
correlation is at least the univariate one by construction, so a
multi-channel model can exceed a univariate ``ceiling'' without any
defect. Cities left of $x = 0$ have a negative univariate ceiling, for
which a percent-of-ceiling figure is undefined.\label{fig:ceiling}}
\end{figure}
\unskip

\subsection{The Thermal-Gap Heuristic Across 49 Cities}

The built--vegetated LST gap (mean LST over impervious cells minus mean
LST over vegetated cells) is the physical driver behind the ceiling, and
is far cheaper to compute. It validates, across the 49-city sample, against
the sound multivariate ceiling, Spearman $\rho = +0.69$
($p = 3 \times 10^{-8}$).

A univariate proxy target gives a more flattering $\rho = +0.84$,
confirmed significantly higher by a paired bootstrap
($\Delta\rho = +0.15$, [$+0.03$, $+0.29$]). We report both because the
difference is diagnostic. Part of the univariate figure's apparent
strength comes from a shared arid sign-flip affecting both variables, not
from gap predictiveness as such.

The heuristic generalizes across arid-exclusion subsets ($p < 10^{-4}$)
with one exception, the humid half ($n = 24$, $\rho = +0.36$,
$p = 0.08$). We leave this open: at this sample size, underpowering and a concentration of the association in arid cities are not
separable, and we prefer to state as such rather than assert the more convenient
reading.

The association is not an artifact of measuring the ceiling at pixel
scale. Recomputing all 49 ceilings at matched receptive field gives
$\rho = +0.716$ ($p = 7 \times 10^{-9}$) against $+0.694$ published; a
paired bootstrap resampling cities together puts the difference at
$+0.022$ with 95\% CI $[-0.085, +0.142]$. The heuristic therefore
predicts recoverable signal rather than merely per-pixel recoverable
signal, which is what makes it usable as a screen for models of any
receptive field. Under matched ceilings, no US city is separability-limited either: Spokane has the lowest ceiling at $+0.335$, the median is $+0.785$, and no ceiling is negative. Furthermore, all five cities that invert under OpenStreetMap limits recover (Boise $-0.331 \rightarrow +0.747$, Albuquerque
$-0.308 \rightarrow +0.716$, Tucson $-0.254 \rightarrow +0.602$, Phoenix
$-0.095 \rightarrow +0.662$, Billings $-0.034 \rightarrow +0.418$).

In arid climates the land-cover--temperature relationship \emph{inverts}
rather than merely weakening. Natural desert vegetation runs hotter than
irrigated, shaded built surfaces, driving the univariate ceiling negative
(Figure~\ref{fig:inversion}).

That arid cities can read cooler than their surroundings is established:
surface urban cool islands are documented across arid and semi-arid
regions, and global analyses find them widespread and trending
\cite{qiu2025longterm}. Liu et al.\ \cite{liu2024identification} further
show, across 104 arid cities, that whether a given city registers as a
heat or a cool island depends on how urban and rural are delimited, with
annual mean differences between methods exceeding 1~$^{\circ}$C. Our
Section~\ref{sec:caveat} result is consistent with theirs and reached
independently on a different sample.

What is new here is not the inversion but its \emph{scale dependence}. The
cool-island literature treats the effect as a property of the city, to be
measured once a definition is fixed. We find instead that it is a property
of the resolution at which land cover is read: every inverting city in our
sample recovers a strongly positive ceiling once the estimator sees
neighborhoods (Section~\ref{sec:pixelscale}), and the arid/non-arid
difference stops being detectable. The practical consequence is that a
city reported as thermally inverted may be perfectly modellable at a
wider receptive field, which the definition-sensitivity framing does not
predict.

\begin{figure}[H]
\includegraphics[width=\textwidth]{fig_inversion_wide.pdf}
\caption{The inversion rendered as surface relief, where height and
color both encode the plotted variable. The built core rises as a
plateau in \emph{both} cities' impervious channel (\textbf{a},\textbf{c};
range 0--1), but the temperature surface rises with it in Austin
(\textbf{b}; 38--48~$^{\circ}$C) and sinks into a basin in Riyadh
(\textbf{d}; 53--57~$^{\circ}$C), where the surrounding desert is hotter
than the city. Nothing in the land-cover input distinguishes the two
cases, which is precisely why a univariate land-cover ceiling reports
that no model can work in the arid one. Single representative cloud-free
dates; Riyadh under the 25~km disc definition.\label{fig:inversion}}
\end{figure}
\unskip

\subsection{Is the Inversion Climatic or Continental?}

The 49-city sample is entirely American, so ``arid'' is confounded with
``American Southwest'' and the inversion could reflect a continental
artifact, such as a difference in construction practice, irrigation policy or
land-use zoning rather than climate. Adding a sixth Southwest city cannot
resolve this, so we therefore repeat the diagnostic on 28 cities across five
continents, under the uniform disc definition described in
Section~2.1.

The inversion is climatic. Eleven cities show negative univariate
ceilings and nine of them lie outside the US (Table~\ref{tab:arid}),
spanning Arabia, the Levant, North Africa, the Atacama, Iberia, Mexico
and central Australia. Arid and
humid/temperate ceilings differ sharply (median $-0.032$ against
$+0.615$; Mann--Whitney $p = 2.4 \times 10^{-4}$), while arid US and arid
non-US cities do \emph{not} differ detectably ($p = 0.075$)---the
comparison an all-US sample cannot make. The gap heuristic replicates on
this independent sample ($\rho = +0.92$ univariate, $+0.57$
multivariate), as does the aridity link (mean LST against ceiling,
$\rho = -0.69$).

\begin{table}[H]
\caption{Every city with a negative multi-day mean univariate ceiling
under the uniform 25~km disc definition, on the five continents. The context
column is the receptive-field-matched ceiling of
Section~\ref{sec:receptive}: no inverting city remains
low.\label{tab:arid}}
\begin{tabularx}{\textwidth}{llCCC}
\toprule
\textbf{City} & \textbf{Region} & \textbf{Univariate} &
\textbf{Context} & \textbf{Gap (}{\boldmath$^{\circ}$}\textbf{C)} \\
\midrule
Amman & Levant & $-0.429$ & $+0.858$ & $-2.52$ \\
Riyadh & Arabia & $-0.342$ & $+0.612$ & $-1.56$ \\
Marrakesh & North Africa & $-0.156$ & $+0.743$ & $-1.18$ \\
Antofagasta & South America & $-0.124$ & $+0.687$ & $-3.38$ \\
Hermosillo & North America & $-0.116$ & $+0.686$ & $-0.99$ \\
Madrid & Europe & $-0.102$ & $+0.468$ & $-0.51$ \\
Billings \textsuperscript{1} & North America & $-0.101$ & $+0.640$ & $-0.61$ \\
Alice Springs & Australia & $-0.096$ & $+0.559$ & $-3.11$ \\
Cairo \textsuperscript{2} & North Africa & $-0.082$ & $+0.912$ & $+1.08$ \\
Yuma \textsuperscript{1} & North America & $-0.067$ & $+0.866$ & $-0.48$ \\
Mendoza & South America & $-0.032$ & $+0.859$ & $-0.57$ \\
\bottomrule
\end{tabularx}

\noindent{\footnotesize{\textsuperscript{1} United States; the remaining
nine lie outside it. \textsuperscript{2} Cairo is the sample's one
exception to the gap mechanism: it inverts univariately despite a
\emph{positive} built--vegetated gap, and its multivariate ceiling
($+0.834$) is among the highest here. The Nile corridor and irrigated
Delta fringe place intensely cultivated land directly against desert
within the same disc, so the two-class built-versus-vegetated contrast the
gap assumes does not describe the city. We flag it rather than exclude
it.}}
\end{table}

\subsubsection{A Robustness Caveat We Report Rather Than Omit}
\label{sec:caveat}

Inversion in the American Southwest is sensitive to how the city is
delimited. Of the five US cities that invert under OpenStreetMap city
limits, four (Phoenix, Tucson, Albuquerque, Boise) do \emph{not} invert
under the disc definition; only Billings does. The hyper-arid non-US cities invert decisively under either definition, with gaps of $-3.38$~$^{\circ}$C for Antofagasta and $-3.11$~$^{\circ}$C for Alice Springs, indicating that the mechanism itself is robust. By contrast, the US examples lie near the sign boundary and therefore should not be interpreted as definition-independent.

Palm Springs is the converse control. It shows no inversion ($+0.255$)
and the sample's largest positive gap ($+6.34$~$^{\circ}$C), in the same
desert as Yuma, because its landscape is irrigated, exhibiting the mechanism
behaving exactly as predicted in a case selected to test it.

\subsubsection{The Inversion Is a Pixel-Scale Effect}
\label{sec:pixelscale}

Applying the receptive-field-matched ceiling of
Section~\ref{sec:receptive} to all 28 cities of the global extension
changes what the inversion means, though not whether it is real.

Every city gains varying amounts, from $+0.019$ (Chicago) to $+0.593$ (Boise), and the size of the gain is predicted almost entirely by how low the per-cell
ceiling was (Spearman $\rho = -0.869$, $p = 2 \times 10^{-9}$) rather
than by aridity directly ($\rho = +0.310$, $p = 0.11$). Arid cities gain
more (median $+0.328$ against $+0.167$, Mann--Whitney $p = 0.005$)
because they are the cities whose per-cell ceilings were low, and not because
aridity acts on the gain through another separate route.

Every inverting city recovers, and all eleven cities in
Table~\ref{tab:arid} reach a context ceiling between $+0.468$ and
$+0.912$. More strikingly, no city anywhere in the 28-city sample is
separability-limited once the ceiling can see neighborhoods: the lowest context
ceiling is Windhoek's $+0.463$ and the median is $+0.779$. The spread
across cities collapses from $0.858$ at pixel scale to $0.526$ with
context, so much of the apparent variation in data readiness is an
artifact of the scale at which readiness was measured.

Consequently the arid/non-arid ceiling difference, which is clear at
pixel scale (median $+0.372$ against $+0.632$, $p = 0.018$), is no longer
detectable once receptive fields are matched ($+0.727$ against $+0.828$,
$p = 0.115$). The thermal-gap heuristic survives this, weakened but
significant: $\rho = +0.571$ against the linear ceiling, $+0.540$
per-cell and $+0.483$ against the context ceiling ($p = 0.009$).

The inversion is therefore a property of pixel-scale land cover rather
than a limit on modeling arid cities. The built core still separates
thermally from the desert around it; what fails is the assumption that
individual built and vegetated cells separate.

\subsection{Channel Screening Without a Simulator}

Because the ceiling depends only on input data, candidate channels can be
screened before any model is built. Adding MODIS NDVI raises the ceiling
predictably across the 49-city sample. The gain correlates with ceiling
height ($\rho = -0.55$, $p = 4.9 \times 10^{-5}$) and with aridity
($\rho = +0.50$, $p = 2.7 \times 10^{-4}$), with 27 of 49 cities robustly
positive and a median gain of $+0.06$.

Phoenix, the irrigated-desert city that motivated adding the channel in
the first place, ranks 4th of 49 ($+0.22$), behind Bakersfield, Spokane
and Boise. The channel is therefore a general, aridity-driven screening
rule rather than a single favorable city's result, which is the
distinction that determines whether such a screen is worth applying
elsewhere.

\subsection{Evapotranspiration Ablation}

Replacing the linear ET term with canonical Penman--Monteith reduces test
$r$ from $+0.755$ to $+0.416$. Leaf-area-index saturation and
vapor-pressure-deficit decoupling, the two mechanisms we isolated,
together explain only 21\% of this regression. Stomatal resistance $r_s$
recovers 61\% at its optimum ($r_s = 3200$~s~m$^{-1}$, roughly
$32\times$ a reference grass value, and not physically defensible), but
even there Penman--Monteith ($r = +0.686$) underperforms simply
\emph{removing} the ET term ($r = +0.758$; Table~\ref{tab:pm}).

No formulation we tested, including one tuned to a physically
indefensible parameter specifically to maximize its own score, recovers
what deleting the term already achieves. This reframes the question from
``which ET formulation is best'' to ``does this linear ET term do useful
work at 500~m at all''. On this evidence it does not.

\begin{table}[H]
\caption{Evapotranspiration formulation ablation. Each configuration is
independently recalibrated, so the comparison is not confounded by
coefficients fitted for a different model.\label{tab:pm}}
\begin{tabularx}{\textwidth}{lC}
\toprule
\textbf{Configuration} & \textbf{Test} {\boldmath$r$} \\
\midrule
No ET term & {\boldmath$+0.758$} \\
Linear ET (reference) & $+0.755$ \\
PM, $r_s = 3200$~s~m$^{-1}$ (optimum, indefensible) & $+0.686$ \\
PM, canonical ($K = 3$, $r_s = 100$~s~m$^{-1}$) & $+0.416$ \\
\bottomrule
\end{tabularx}
\end{table}

\subsection{Counterfactual Scenarios}

Under ship-configuration ERA5 forcing, a $+20\%$ canopy increase produces
a mean $\Delta T = -0.74$~$^{\circ}$C, a citywide slope of
$-0.44$~$^{\circ}$C per 10\% canopy, within the literature band
\cite{zaerpour2025tree,wri2026cooling}. A river greenway yields
$-0.25$~$^{\circ}$C and cool roofs $-0.07$~$^{\circ}$C
\cite{santamouris2014cool,wang2022berlin}; suburban densification, run as
a control in the opposite direction, yields $+0.18$~$^{\circ}$C.

These means are averaged over every (cell, timestep) pair in the 48-hour
run. The convention matters: the same scenario gives
$-1.52$~$^{\circ}$C at the final frame alone, and per-frame values swing
between roughly $-0.3$ and $-1.5$~$^{\circ}$C across the diurnal cycle,
so a cooling figure quoted without its averaging window is ambiguous by a
factor of two. The spatial pattern is likewise non-uniform: cooling
concentrates where there was impervious surface to convert, and water
cells warm slightly relative to baseline.

Two quantities fall outside their literature bands. Local per-patch
cooling ($-0.59$~$^{\circ}$C per 10\%) is 1.5--5$\times$ below field
measurements \cite{krayenhoff2021albedo}, and the edge-referenced
half-cooling distance ($\approx$330~m) sits just outside the 100--300~m
band. Both reflect the diffusion trade-off analyzed below rather than
independent defects.

\section{Discussion}
\label{sec:disc}

\subsection{What Became of the Data Limit}
\label{sec:dichotomy}

This paper began from a dichotomy: a low transfer score means either that
the model failed or that the input data could not support any model, and the
diagnostic was built to tell those apart. It is worth stating plainly what
happened when we applied it, because the result is not the one that the framing
anticipates.

Across 49 US cities under administrative boundaries and 28 cities under a
uniform disc definition, \emph{no city is separability-limited} once the ceiling
is specified to match the model's receptive field. The lowest matched
ceiling anywhere is $+0.335$; the median is near $+0.78$ in both samples.
Every city-definition whose univariate ceiling is negative (sixteen of
them, eleven in the global sample and five in the US, the cases that
motivated the arid-inversion result) recovers to a strongly positive
matched ceiling. The second branch of the dichotomy, as measured by land
cover, is close to unoccupied.

We read this as a positive finding instead of a null one. Low
land-cover-versus-temperature correlations are real and are routinely
interpreted as evidence that a city is intrinsically hard to model. Our
results say that interpretation is usually a specification artifact: the
correlation was measured at one scale and the model operates at another.
The practical question shifts from \emph{whether} land cover can support a
model in a given city to \emph{at what scale} it does, which is a question
with an actionable answer everywhere we looked.

\subsection{The Real Data Limit Is Observability, Not Separability}

However, the dichotomy is not empty. It is occupied by a failure mode the
ceiling does not measure.

Of 35 cities attempted in the global extension, seven yielded no usable
warm-season observation at all, not having a weak ceiling, but zero cloud-free
days clearing our coverage threshold. The US sweep lost a further six of
55 the same way. These cities have no ceiling because they have no data,
so the diagnostic is silent on them by our construction, since it is computed
from scenes that exist.

The losses are not confined to one climate. Bangkok, Manila and Lagos are
humid tropical, but Dubai, Kuwait City and Jodhpur are arid, where dust
and haze trigger the same cloud masking as cloud does. Roughly one city in
seven across both samples was lost this way, and the loss falls hardest on the places where heat exposure is most severe and where modeling
capacity is scarcest.

This also conditions our own headline. That no city is
separability-limited is measured among the cities we could observe, and
observability is not random with respect to climate: a city with no
usable scene is limited by its data in a way $C(X,Y;R)$ cannot register,
because $C(X,Y;R)$ is computed from scenes that exist. A complete data-readiness assessment therefore
needs two axes, not one: whether land cover separates thermal regimes at
some accessible scale, and whether the thermal observations needed to
check that exist in the first place. We report the first here and can only
flag the second, which we regard as the more urgent gap for any attempt to
extend this diagnostic globally.

\subsection{What the Physics Results Are For}

The two results that follow concern the simulator rather than the
diagnostic, and are included for one reason: a ceiling diagnostic is only
as informative as the model it is demonstrated on. A reader is entitled to
ask whether the transfer failures in Section~\ref{sec:rfresult} reflect
the diagnostic or merely a weak instrument, and these establish what the
instrument can and cannot do. They are not offered as independent
contributions.

\subsection{The Diffusion--Fidelity Trade-Off}

No single Fickian diffusivity $D$ satisfies both MODIS spatial
autocorrelation and literature cooling-decay scales. The best $r$
achievable inside the 100--300~m band is $+0.67$ at $D = 0.20$, a cost of
$-0.19$ in $r$ against the unconstrained optimum ($r = +0.85$ at
$D \approx 12.6$). This is roughly twice the cost a naive, uncalibrated
estimate would suggest, and it means a model tuned purely for pattern
correlation will produce intervention footprints too diffuse to use for
planning.

\subsection{Advection Is Not a Drop-In Fix}

Replacing diffusion with ERA5 wind-driven advection destroys skill
($r = -0.11$ at 25\% strength). The reason is structural rather than a
matter of tuning: a spatially uniform wind translates the entire surface
pattern downwind, whereas the physical process transports heat in the air
layer above the surface. A coupled air layer, not a one-line
substitution, is required \cite{mitraka2025prototypes}.

\subsection{Class Separability as an A Priori Screen}

Because the ceiling depends only on input data, it can be computed for a
candidate city before any modeling effort is committed, and this is where
the diagnostic earns its keep for practitioners rather than for
methodologists. A city considering heat-mitigation modeling faces a
concrete decision: whether to invest in calibrating and deploying a model
at all, and if so at what resolution its land-cover inventory needs to be
read. Both questions are answerable in advance, from a boundary, a
land-cover raster and one cloud-free scene, at a cost measured in minutes
rather than in a modeling programme.

Given Section~\ref{sec:dichotomy}, what the screen buys is not a shortlist
of cities worth modeling---on present evidence that is most of them---but
the scale at which each city's land cover becomes actionable. That is the
operationally useful output. A municipality whose land cover only becomes
informative at 1--2~km should not expect a pixel-scale model to guide
street-level planting decisions, and will learn this before commissioning
one rather than after. For a deployment across many cities the ordering
reverses accordingly: establish the informative scale first, then
calibrate at that scale, rather than discovering the mismatch through a
poor transfer score and misreading it as an absence of signal.

What the screen does \emph{not} do is diagnose why a given model fails.
Our own results make the distinction concrete. Under per-cell ceilings
Phoenix looks separability-limited and Denver looks recoverable; under
receptive-field-matched ceilings both have ample headroom
(Section~\ref{sec:rfresult}), and Phoenix's poor transfer becomes a
statement about this simulator rather than about Phoenix. A low ceiling
is evidence that no model of the matched class will do well. A high
ceiling with a low score localizes the problem to the model, which is
useful, but it does not say which part.

\subsection{Scale, Not Aridity, Is What Limits These Cities}

Section~\ref{sec:pixelscale} changes the practical reading of the arid
result. Arid cities are not intrinsically resistant to land-cover-driven
modeling; they are poorly served by \emph{pixel-scale} land-cover models
specifically. The remedy indicated is a wider receptive field rather than
additional channels, and the two are not interchangeable: widening the
receptive field raises Phoenix's ceiling by $+0.31$, whereas adding an
NDVI channel raises it by $+0.22$ at pixel scale and leaves the
pixel-scale limitation in place.

It is worth being precise about what aridity does here, because the
obvious reading is wrong. Every city in both samples gains from a wider
receptive field (all 49 US cities, minimum $+0.026$), so this is not a
correction that applies to arid cities alone. What varies is the size of
the gain, which tracks how low the per-cell ceiling was (Spearman
$\rho = -0.707$ in the US sample, $-0.869$ globally) far more closely than
it tracks temperature ($+0.208$, $p = 0.15$; $+0.310$, $p = 0.11$).
Splitting the US sample at median LST, hot and cool halves gain
indistinguishably ($+0.212$ against $+0.225$, $p = 0.73$). Arid cities
gain more in the global sample because their per-cell
ceilings are low, not because aridity acts through some separate route.

This reframes what a readiness assessment should report, as a single
per-city number presumes the ceiling is a property of the city. It is
better understood as a property of the city \emph{and} the scale at which
its land cover is read, and the useful output for a practitioner is the
scale at which land cover becomes informative there, which is wider than
one pixel essentially everywhere, and widest where a pixel-scale
assessment would have declared the city hopeless.

\subsection{Structure in the Residual Bounds What Calibration Can Buy}

The residual in Figure~\ref{fig:validation}f is spatially structured
rather than noise-like. The model reproduces the hot core but places it
slightly west of where MODIS observes it, leaving a coherent dipole. This
is a positional error, not random misfit, and it implies that
$r = +0.75$ is limited by something further coefficient tuning cannot
remove. Three candidate causes, in increasing cost to test, are the
diffusion trade-off smearing the core, an absorption composite too coarse
to place it precisely, and a peak-frame selection offset from the Aqua
overpass time. Establishing which dominates would bound what calibration
can still achieve, and we flag it as the most informative next diagnostic.

\subsection{Limitations}

The simplified longwave scheme leaves a day-dependent absolute offset
rather than a constant bias. Across the eight Austin focus days the mean
simulated-minus-observed offset is $-1.40$~$^{\circ}$C, ranging from
$-5.47$ to $+3.06$~$^{\circ}$C. It couples only weakly to spatial skill
($\Delta r = 0.02$ across a 3~$^{\circ}$C bias sweep), but a swing of
8.5~$^{\circ}$C between days means absolute temperatures should not be
read as calibrated predictions. The 500~m grid smooths
highway-corridor hot streaks that finer imagery resolves, and per-city
ceilings should be read as multi-day estimates, since single-day values
can be actively unrepresentative.

Two results remain open at the present sample sizes: the humid-subset
association of Section~3.4, and the definition-sensitivity of the US arid
cities.

A third limitation is one of specification. Our matched ceiling uses
three Gaussian scales on three channels, chosen to bracket the shipped
diffusion length, rather than a search over receptive fields. The matched
ceiling is therefore at least as high as we report and may be higher,
which makes the direction of every conclusion here safe and the
magnitudes lower bounds. A model class with a wider or differently shaped
receptive field would need its own matched ceiling before a
percent-of-ceiling figure against it would mean anything---a six-layer
convolutional network already exceeds our specification in every focus
city (Section~\ref{sec:cnn}), so this is a demonstrated property of the
framework rather than a hypothetical caveat.

\section{Conclusions}

A land-cover data ceiling separates two failure modes that a transfer
score alone conflates, but only when it is specified to match the model it
is meant to bound. The quantity that does the separating is not an upper
bound on prediction. It is an achievable-skill benchmark conditioned on a
stated information set and a stated spatial support, $C(X,Y;R)$, and its
usefulness comes from that conditioning rather than in spite of it: a
diagnostic computed at the wrong support does not merely lose precision,
it reports a data limitation where the limitation is the model's. Firstly, the ceiling must span the channels a model
consumes, as the univariate alternative goes negative in arid cities, where
a percent-of-ceiling figure becomes undefined, and its receptive field must
match the model's, since a ceiling regressed on each cell's own land cover
does not bound a model that reads across space. This distinction is
consequential in practice: our simulator exceeds the per-cell ceiling on
the days used for evaluation in Austin, demonstrating that a mismatch in
receptive field can invalidate the intended bound.

Five results transfer beyond the simulator used to demonstrate them.
First, the ceiling diagnostic itself, validated across 49 cities and made
actionable through a cheap thermal-gap heuristic and a predictable NDVI
channel screen. Second, evidence that the arid inversion is climatic
rather than continental, established on 28 cities across five continents
under a uniform city definition, while being explicit that the American
Southwest cases are definition-sensitive.

Third, and the result we did not expect, is that the inversion is a
pixel-scale effect. Across both samples every inverting city recovers
once the ceiling is allowed to read neighborhoods, to between $+0.47$
and $+0.91$ globally, and $+0.42$ to $+0.75$ in the US, and no observed city in either sample
remained separability-limited under the
matched specification (global minimum $+0.463$,
US minimum
$+0.335$). The arid/non-arid ceiling difference ceases to be detectable
globally ($p = 0.018$ at pixel scale, $p = 0.115$ matched). What limits
these cities is the scale at which land cover is read, not the city, and
the correction applies everywhere rather than to arid cities alone: all
49 US cities gain, with the size of the gain tracking how low the
per-cell ceiling was rather than how hot the city is. A readiness
assessment should therefore report the scale at which land cover becomes
informative rather than a single per-city number.

Fourth, a diffusion--fidelity trade-off roughly twice as steep as a naive
estimate. Fifth, an evapotranspiration ablation showing that a standard
linear ET term contributes nothing to spatial-pattern skill at 500~m,
regardless of formulation.

The simulator itself achieves $r = +0.754$ on held-out MODIS days using
only inputs retrievable by anonymous HTTP request, which makes every
result above independently checkable without an institutional data
agreement.

\vspace{6pt}

\funding{This research received no external funding.}

\institutionalreview{Not applicable.}

\informedconsent{Not applicable.}

\dataavailability{All input data are publicly available without
authentication: MODIS MYD11A1 and MOD13A1 via the Microsoft Planetary
Computer STAC API, ESA WorldCover v200 via the Copernicus/ESA object
store, ERA5 via the Google ARCO-ERA5 Zarr mirror, and administrative
boundaries via OpenStreetMap. Source code, calibration artifacts and the
scripts that regenerate every figure and table reported here are
available at \url{https://github.com/abhi-is-cool/austin-twin}.}

% REVIEW: the GenAI sentence below must describe what you actually did.
\acknowledgments{MODIS data are provided via the Microsoft Planetary
Computer; ESA WorldCover via Copernicus/ESA; ERA5 via Google ARCO-ERA5 on
Google Cloud Storage; OpenStreetMap data are copyright OpenStreetMap
contributors and licensed under the ODbL. During the preparation of this study, the author used 
artificial intelligence for software implementation, analysis scripting, and figure generation.
All research questions, methodological decisions and interpretations are the author's own.
The author has reviewed and edited the output and takes full responsibility for the content of this
publication.}

\conflictsofinterest{The author declares no conflicts of interest.}

\abbreviations{Abbreviations}{%
The following abbreviations are used in this manuscript:\\

\noindent
\begin{tabular}{@{}ll}
ARCO & Analysis-Ready, Cloud-Optimized\\
CFL & Courant--Friedrichs--Lewy\\
ET & Evapotranspiration\\
LAI & Leaf Area Index\\
LST & Land Surface Temperature\\
NDVI & Normalized Difference Vegetation Index\\
PM & Penman--Monteith\\
STAC & SpatioTemporal Asset Catalog\\
UHI & Urban Heat Island\\
UTM & Universal Transverse Mercator\\
VPD & Vapor Pressure Deficit
\end{tabular}
}

\begin{adjustwidth}{-\extralength}{0cm}

\reftitle{References}

\isAPAandChicago{}{%
\begin{thebibliography}{999}

\bibitem[Bhamjee et al.(2024)]{bhamjee2024detection}
Bhamjee, M.; Debary, H.; Gaffoor, Z.; Govindasamy, T.R.; Mahlasi, C.;
Fiaz, M.; Vos, E.E.; Klein, L.; Makhanya, S.; Watson, C.; Kuehnert, J.
Detection and characterization of urban heat islands with machine
learning. In Proceedings of the 2024 IEEE International Geoscience and
Remote Sensing Symposium (IGARSS 2024), Athens, Greece, 7--12 July 2024;
pp. 1693--1699. \url{https://doi.org/10.1109/IGARSS53475.2024.10641750}.

\bibitem[Bueno et al.(2013)]{bueno2013urban}
Bueno, B.; Norford, L.; Hidalgo, J.; Pigeon, G. The urban weather
generator. \textit{J. Build. Perform. Simul.} \textbf{2013}, \textit{6},
269--281. \url{https://doi.org/10.1080/19401493.2012.718797}.

\bibitem[Cover and Thomas(2006)]{cover2006elements}
Cover, T.M.; Thomas, J.A. \textit{Elements of Information Theory}, 2nd ed.;
Wiley-Interscience: Hoboken, NJ, USA, 2006.

\bibitem[Noshad et al.(2019)]{noshad2019learning}
Noshad, M.; Xu, L.; Hero, A. Learning to benchmark: Determining best
achievable misclassification error from training data. \textit{arXiv}
\textbf{2019}, arXiv:1909.07192.
\url{https://doi.org/10.48550/arXiv.1909.07192}.

\bibitem[Friedman(2001)]{friedman2001greedy}
Friedman, J.H. Greedy function approximation: A gradient boosting machine.
\textit{Ann. Stat.} \textbf{2001}, \textit{29}, 1189--1232.
\url{https://doi.org/10.1214/aos/1013203451}.

\bibitem[Hersbach et al.(2020)]{hersbach2020era5}
Hersbach, H.; Bell, B.; Berrisford, P.; Hirahara, S.; Hor\'anyi, A.;
Mu\~noz-Sabater, J.; Nicolas, J.; Peubey, C.; Radu, R.; Schepers, D.;
Simmons, A.; Soci, C.; Abdalla, S.; Abellan, X.; Balsamo, G.;
Bechtold, P.; Biavati, G.; Bidlot, J.; Bonavita, M.; De Chiara, G.;
Dahlgren, P.; Dee, D.; Diamantakis, M.; Dragani, R.; Flemming, J.;
Forbes, R.; Fuentes, M.; Geer, A.; Haimberger, L.; Healy, S.;
Hogan, R.J.; H\'olm, E.; Janiskov\'a, M.; Keeley, S.; Laloyaux, P.;
Lopez, P.; Lupu, C.; Radnoti, G.; de Rosnay, P.; Rozum, I.; Vamborg, F.;
Villaume, S.; Th\'epaut, J.-N. The ERA5 global reanalysis.
\textit{Q. J. R. Meteorol. Soc.} \textbf{2020}, \textit{146}, 1999--2049.
\url{https://doi.org/10.1002/qj.3803}.

\bibitem[Hoang and Nguyen(2025)]{hoang2025geospatial}
Hoang, N.D.; Nguyen, Q.L. Geospatial analysis and machine learning
framework for urban heat island intensity prediction: Natural gradient
boosting and deep neural network regressors with multisource remote
sensing data. \textit{Sustainability} \textbf{2025}, \textit{17}, 4287.
\url{https://doi.org/10.3390/su17104287}.

\bibitem[Reinbigler et al.(2026)]{heatmat2026}
Reinbigler, M.; Rouffet, R.; Naylor, P.; Czerkawski, M.; Dionelis, N.;
Brunet, E.; Fetita, C.; Martin, R. HeatMat: Simulation of city material
impact on urban heat island effect. \textit{Comput. Graph. Forum}
\textbf{2026}. \url{https://doi.org/10.1111/cgf.70376}.

\bibitem[Krayenhoff et al.(2021)]{krayenhoff2021albedo}
Krayenhoff, E.S.; Broadbent, A.M.; Zhao, L.; Georgescu, M.; Middel, A.;
Voogt, J.A.; Martilli, A.; Sailor, D.J.; Erell, E. Cooling hot cities: A
systematic and critical review of the numerical modelling literature.
\textit{Environ. Res. Lett.} \textbf{2021}, \textit{16}, 053007.
\url{https://doi.org/10.1088/1748-9326/abdcf1}.

\bibitem[Kusaka et al.(2001)]{kusaka2001simple}
Kusaka, H.; Kondo, H.; Kikegawa, Y.; Kimura, F. A simple single-layer
urban canopy model for atmospheric models: Comparison with multi-layer
and slab models. \textit{Bound.-Layer Meteorol.} \textbf{2001},
\textit{101}, 329--358. \url{https://doi.org/10.1023/A:1019207923078}.

\bibitem[Liu et al.(2024)]{liu2024identification}
Liu, Z.; Ye, R.; Yang, Q.; Hu, T.; Liu, Y.; Chakraborty, T.C.; Liao, Z.
Identification of surface urban heat versus cool islands for arid cities
depends on the choice of urban and rural definitions.
\textit{Sci. Total Environ.} \textbf{2024}, \textit{951}, 175631.
\url{https://doi.org/10.1016/j.scitotenv.2024.175631}.

\bibitem[Lyu et al.(2022)]{lyu2022integrated}
Lyu, F.; Wang, S.; Han, S.Y.; Catlett, C.; Wang, S. An integrated cyberGIS
and machine learning framework for fine-scale prediction of urban heat
island using satellite remote sensing and urban sensor network data.
\textit{Urban Inform.} \textbf{2022}, \textit{1}, 6.
\url{https://doi.org/10.1007/s44212-022-00002-4}.

\bibitem[Masson(2000)]{masson2000physically}
Masson, V. A physically-based scheme for the urban energy budget in
atmospheric models. \textit{Bound.-Layer Meteorol.} \textbf{2000},
\textit{94}, 357--397. \url{https://doi.org/10.1023/A:1002463829265}.

\bibitem[Microsoft(2023)]{microsoft_pc}
Microsoft. Planetary Computer STAC API. Available online:
\url{https://planetarycomputer.microsoft.com} (accessed on 1 September
2026).

\bibitem[Mitraka et al.(2025)]{mitraka2025prototypes}
Mitraka, Z.; Lantzanakis, G.; Panagiotakis, E.; Chrysoulakis, N.;
Somarakis, G.; Marconcini, M.; Feigenwinter, C.; Lauwaet, D. Prototypes
for enhanced urban heat monitoring leveraging Copernicus services and
satellite data. \textit{Discov. Cities} \textbf{2025}, \textit{2}, 14.
\url{https://doi.org/10.1007/s44327-025-00059-1}.

\bibitem[Oke(1982)]{oke1982energetic}
Oke, T.R. The energetic basis of the urban heat island. \textit{Q. J. R.
Meteorol. Soc.} \textbf{1982}, \textit{108}, 1--24.
\url{https://doi.org/10.1002/qj.49710845502}.

\bibitem[Qiu et al.(2025)]{qiu2025longterm}
Qiu, R.; Zhan, Q.; Tu, S. Long-term global trends and influencing factors
of surface urban cool and heat islands. \textit{Int. J. Digit. Earth}
\textbf{2025}, \textit{18}, 2489731.
\url{https://doi.org/10.1080/17538947.2025.2489731}.

\bibitem[Santamouris(2014)]{santamouris2014cool}
Santamouris, M. Cooling the cities---A review of reflective and green
roof mitigation technologies to fight heat island and improve comfort in
urban environments. \textit{Sol. Energy} \textbf{2014}, \textit{103},
682--703. \url{https://doi.org/10.1016/j.solener.2012.07.003}.

\bibitem[Storn and Price(1997)]{storn1997differential}
Storn, R.; Price, K. Differential evolution---A simple and efficient
heuristic for global optimization over continuous spaces.
\textit{J. Glob. Optim.} \textbf{1997}, \textit{11}, 341--359.
\url{https://doi.org/10.1023/A:1008202821328}.

\bibitem[Voogt and Oke(2003)]{voogt2003thermal}
Voogt, J.A.; Oke, T.R. Thermal remote sensing of urban climates.
\textit{Remote Sens. Environ.} \textbf{2003}, \textit{86}, 370--384.
\url{https://doi.org/10.1016/S0034-4257(03)00079-8}.

\bibitem[Wang et al.(2022)]{wang2022berlin}
Wang, X.; Li, H.; Sodoudi, S. The effectiveness of cool and green roofs in
mitigating urban heat island and improving human thermal comfort.
\textit{Build. Environ.} \textbf{2022}, \textit{217}, 109082.
\url{https://doi.org/10.1016/j.buildenv.2022.109082}.

\bibitem[Wesley et al.(2026)]{wri2026cooling}
Wesley, E.J.; Mackres, E.; Pool, J.R. Cooling Potential of Urban Trees.
World Resources Institute, 17 April 2026. Available online:
\url{https://www.wri.org/insights/urban-trees-cooling-potential}
(accessed on 21 September 2026).

\bibitem[Zanaga et al.(2022)]{zanaga2022esa}
Zanaga, D.; Van De Kerchove, R.; Daems, D.; De Keersmaecker, W.;
Brockmann, C.; Kirches, G.; Wevers, J.; Cartus, O.; Santoro, M.;
Fritz, S.; Lesiv, M.; Herold, M.; Tsendbazar, N.E.; Xu, P.; Ramoino, F.;
Arino, O. ESA WorldCover 10~m 2021 v200. \textit{Zenodo} \textbf{2022}.
\url{https://doi.org/10.5281/zenodo.7254221}.

\bibitem[Zaerpour et al.(2025)]{zaerpour2025tree}
Zaerpour, M.; Papalexiou, S.M.; Pietroniro, A. Increasing tree canopy
lowers urban air temperature by up to 1.5~$^{\circ}$C in heat-prone areas.
\textit{npj Urban Sustain.} \textbf{2025}, \textit{5}, 92.
\url{https://doi.org/10.1038/s42949-025-00277-x}.

\bibitem[Ziter et al.(2019)]{ziter2019scale}
Ziter, C.D.; Pedersen, E.J.; Kucharik, C.J.; Turner, M.G. Scale-dependent
interactions between tree canopy cover and impervious surfaces reduce
daytime urban heat during summer. \textit{Proc. Natl. Acad. Sci. USA}
\textbf{2019}, \textit{116}, 7575--7580.
\url{https://doi.org/10.1073/pnas.1817561116}.

\end{thebibliography}
}

\PublishersNote{}
\end{adjustwidth}
\end{document}
